\chapter{Methodes et architectures}

\section{Idee generale que je retiens de PointNet}
La forme que je retiens de l'article est:
\begin{align}
f(\{x_1, \dots, x_n\}) \approx \gamma\left(\max_i\, h(x_i)\right)
\end{align}
avec \(h\) un encodeur partage par point et \(\max\) une aggregation symetrique.
Cette ecriture m'a servi de fil conducteur pour comprendre pourquoi PointNet fonctionne mieux qu'un MLP classique sur des donnees non ordonnees.

\section{Reseau 1: PointMLP (Exercice 1)}
J'ai d'abord implemente le modele demande dans l'enonce: flatten \(1024\times 3\) puis couches fully-connected.

\begin{figure}[H]
\centering
\begin{tikzpicture}[node distance=0.75cm]
    \node[input] (in) {Input Point Cloud \\ $(N \times 3)$};
    \node[layer, below=of in] (flatten) {Flatten \\ $(3072)$};
    \node[layer, below=of flatten] (fc1) {FC + BN + ReLU \\ $(3072 \rightarrow 512)$};
    \node[layer, below=of fc1] (fc2) {FC + BN + ReLU \\ $(512 \rightarrow 256)$};
    \node[layer, below=of fc2] (drop) {Dropout \\ $(p=0.3)$};
    \node[layer, below=of drop] (fc3) {FC \\ $(256 \rightarrow K)$};
    \node[layer, below=of fc3] (logsoft) {LogSoftmax};
    \node[input, fill=red!10, below=of logsoft] (out) {Output Classes \\ $(K)$};

    \draw[arrow] (in) -- (flatten);
    \draw[arrow] (flatten) -- (fc1);
    \draw[arrow] (fc1) -- (fc2);
    \draw[arrow] (fc2) -- (drop);
    \draw[arrow] (drop) -- (fc3);
    \draw[arrow] (fc3) -- (logsoft);
    \draw[arrow] (logsoft) -- (out);
\end{tikzpicture}
\caption{Architecture du PointMLP (conforme au notebook: 3072\(\rightarrow\)512\(\rightarrow\)256\(\rightarrow\)K).}
\label{fig:arch-mlp}
\end{figure}

\noindent\textbf{Analyse.} Avec ce modele, je perds la structure "ensemble de points" au moment du flatten.
Le modele voit un grand vecteur ordonne, alors que l'ordre des points est arbitraire.
C'est la limite qui m'a pousse a passer a PointNetBasic.

\section{Reseau 2: PointNetBasic (Exercice 2.1)}
Dans PointNetBasic, je remplace le flatten par des MLP partages (\texttt{Conv1d} de noyau 1) puis un max-pooling global.

\begin{figure}[H]
\centering
\resizebox{0.98\linewidth}{!}{%
\begin{tikzpicture}[node distance=0.6cm and 1cm]
    \node[input] (in) {Input Cloud \\ $(B, 3, N)$};

    \node[layer, right=of in] (conv1) {Conv1D(64) + BN + ReLU};
    \node[layer, right=of conv1] (conv2) {Conv1D(64) + BN + ReLU};
    \node[layer, right=of conv2] (conv3) {Conv1D(64) + BN + ReLU};

    \node[layer, below=of conv3] (conv4) {Conv1D(128) + BN + ReLU};
    \node[layer, left=of conv4] (conv5) {Conv1D(1024) + BN + ReLU};

    \node[pool, left=of conv5] (pool) {Global Max Pooling \\ $(B, 1024, 1)$};
    \node[layer, below=of pool] (fc1) {FC(512) + BN + ReLU};
    \node[layer, right=of fc1] (fc2) {FC(256) + BN + ReLU};
    \node[layer, right=of fc2] (drop) {Dropout};
    \node[layer, right=of drop] (fc3) {FC(K) + LogSoftmax};

    \node[input, fill=red!10, below=of fc3] (out) {Output $(B, K)$};

    \draw[arrow] (in) -- (conv1);
    \draw[arrow] (conv1) -- (conv2);
    \draw[arrow] (conv2) -- (conv3);
    \draw[arrow] (conv3) -- (conv4);
    \draw[arrow] (conv4) -- (conv5);
    \draw[arrow] (conv5) -- (pool);
    \draw[arrow] (pool) -- (fc1);
    \draw[arrow] (fc1) -- (fc2);
    \draw[arrow] (fc2) -- (drop);
    \draw[arrow] (drop) -- (fc3);
    \draw[arrow] (fc3) -- (out);
\end{tikzpicture}
}
\caption{Architecture du PointNetBasic (sans T-Net), basee sur les blocs \texttt{Conv1d} du notebook.}
\label{fig:arch-basic}
\end{figure}

\noindent\textbf{Analyse.} Ce schema corrige le probleme principal de PointMLP:
le max-pooling rend la sortie invariante a la permutation des points.
Pour le traitement de point clouds, c'est une difference fondamentale.

\section{Reseau 3: PointNetFull avec T-Net 3x3 (Exercice 2.2)}
La version Full ajoute un T-Net d'alignement 3x3 en entree.
J'utilise aussi une regularisation orthogonale:
\begin{align}
\mathcal{L} = \mathcal{L}_{CE} + \lambda\lVert I-AA^\top\rVert_F^2
\end{align}

\begin{figure}[H]
\centering
\resizebox{0.98\linewidth}{!}{%
\begin{tikzpicture}[node distance=0.8cm and 1.5cm]
    \node[input] (in) {Input Cloud \\ $(B, N, 3)$};

    \node[tnet, below right=of in, yshift=1cm] (tnet) {Tnet ($3 \times 3$)};
    \node[op, right=of in] (mul) {$\otimes$};

    \node[layer, right=of mul] (features) {PointNetBasic \\ Feature Extraction \\ (Conv1D jusqu'a MaxPool)};
    \node[layer, right=of features] (mlp) {MLP Head \\ (FC 512 \(\rightarrow\) 256 \(\rightarrow\) K)};
    \node[input, fill=red!10, right=of mlp] (out) {Scores \\ $(B, K)$};

    \draw[arrow] (in) -- (mul);
    \draw[arrow] (in) |- (tnet);
    \draw[arrow] (tnet) -| node[pos=0.25, right] {Matrice de transformation ($3 \times 3$)} (mul);
    \draw[arrow] (mul) -- node[above] {$(B, 3, N)$} (features);
    \draw[arrow] (features) -- node[above] {Global Feature $(B, 1024)$} (mlp);
    \draw[arrow] (mlp) -- (out);
\end{tikzpicture}
}
\caption{Architecture du PointNetFull avec module \texttt{Tnet} d'alignement en entree.}
\label{fig:arch-full}
\end{figure}

\noindent\textbf{Analyse.} Le T-Net apporte une adaptation geometrique supplementaire.
Dans ma lecture, il agit comme une normalisation apprenable: le reseau apprend a "mettre en forme" le nuage avant extraction des features globales.

\section{Comparaison structurelle des trois reseaux}
\begin{table}[H]
\centering
\caption{Comparaison de mes trois implementations.}
\label{tab:archi-compare}
\begin{tabular}{lccc}
\toprule
Modele & Invariance ordre & Alignement geometrique & Params trainables \\
\midrule
PointMLP & Non & Non & 1,708,810 \\
PointNetBasic & Oui (max-pooling) & Non & 811,914 \\
PointNetFull & Oui (max-pooling) & Oui (T-Net 3x3) & 1,614,995 \\
\bottomrule
\end{tabular}
\end{table}

\noindent\textbf{Analyse.} J'insiste sur un point: l'amelioration ne vient pas simplement d'avoir plus de parametres.
Elle vient surtout du fait que PointNetBasic et PointNetFull respectent mieux la nature d'un point cloud (ensemble non ordonne, variations geometriques).
